% A Testing Harness for Local LLM Evaluation in Neuro-Mimetic Systems
% Michael Clark — Are-Self Research
\documentclass[man,12pt,floatsintext]{apa7}

\usepackage[utf8]{inputenc}
\usepackage[T1]{fontenc}
\usepackage{amsmath,amssymb,amsfonts}
\usepackage{graphicx}
\usepackage{booktabs}
\usepackage{hyperref}
\usepackage{algorithm}
\usepackage{algpseudocode}
\usepackage{listings}
\usepackage{xcolor}

\definecolor{codebg}{rgb}{0.95,0.95,0.97}
\definecolor{codegreen}{rgb}{0.2,0.6,0.2}
\definecolor{codegray}{rgb}{0.5,0.5,0.5}
\definecolor{codepurple}{rgb}{0.58,0,0.82}

\lstset{
  backgroundcolor=\color{codebg},
  basicstyle=\ttfamily\small,
  breaklines=true,
  captionpos=b,
  commentstyle=\color{codegreen},
  keywordstyle=\color{codepurple},
  numberstyle=\tiny\color{codegray},
  stringstyle=\color{codegreen},
  frame=single,
  numbers=left,
}

\usepackage[style=apa,backend=biber]{biblatex}
\addbibresource{../../templates/references.bib}

\title{A Testing Harness for Local LLM Evaluation in Neuro-Mimetic Systems}
\shorttitle{LLM TESTING HARNESS}
\author{Michael Clark}
\affiliation{Independent Researcher \\ scipraxian}

\abstract{
Evaluating local language models for autonomous agent workflows presents unique challenges:
models must be tested not just for generation quality but for tool-calling reliability,
format compliance, context window utilization, and behavior under the resource constraints
of consumer hardware. Standard LLM benchmarks (MMLU, HumanEval, etc.) measure general
capability but do not predict whether a model will function correctly as an autonomous
agent in a structured orchestration system. We describe a testing harness built within
Are-Self's architecture that evaluates models against agent-specific criteria: tool call
format compliance, multi-turn consistency, focus budget efficiency, memory formation
quality, and graceful degradation under context pressure. The harness leverages Are-Self's
existing execution infrastructure---spike trains, reasoning sessions, and telemetry---to
run evaluation suites as standard pathways, enabling continuous model assessment as new
quantizations and fine-tunes become available via Ollama.
}

\keywords{LLM evaluation, local inference, tool calling, autonomous agents, model testing,
Ollama, quantization effects}

\begin{document}
\maketitle

\section{Introduction}

The local LLM ecosystem evolves rapidly. New models, quantizations, and fine-tunes appear
weekly on Ollama, Hugging Face, and similar platforms. For systems like Are-Self that
depend on models functioning as autonomous agents---calling tools, managing context,
forming memories---each new model requires evaluation along dimensions that standard
benchmarks do not measure.

A model that scores well on MMLU may hallucinate tool names. A model that excels at
HumanEval may fail to format tool calls in the JSON schema expected by the orchestrator.
A model that handles single-turn generation beautifully may degenerate across 50 turns
of multi-turn reasoning.

We present a testing harness that evaluates models within the actual orchestration
environment, measuring the capabilities that matter for autonomous agent operation.

\section{Related Work}

\subsection{Standard LLM Benchmarks}

% TODO: Survey of MMLU, HumanEval, MATH, BigBench
% - What they measure
% - Why they don't predict agent performance

\subsection{Agent-Specific Evaluation}

% TODO: Survey of AgentBench, ToolBench, MINT
% - Closer to our needs but designed for cloud models
% - Don't account for quantization artifacts

\subsection{Local Model Evaluation}

% TODO: Existing approaches to evaluating quantized/local models
% - Perplexity comparisons across quantizations
% - Task-specific leaderboards
% - Gap: no agent-workflow-specific evaluation for local models

\section{Evaluation Dimensions}

\subsection{Tool Call Compliance}

Does the model emit tool calls in the expected format? We test:

\begin{enumerate}
  \item Correct JSON structure for function calls
  \item Tool name accuracy (no hallucinated tool names)
  \item Parameter type compliance (string vs. integer vs. boolean)
  \item Required parameter inclusion
  \item Graceful handling of tool results
\end{enumerate}

Are-Self's Parietal Lobe (hallucination armor) already validates tool calls; the harness
measures the rate at which calls pass validation.

\subsection{Multi-Turn Consistency}

Over $N$ turns, does the model:

\begin{enumerate}
  \item Maintain coherent task focus
  \item Reference previous turns accurately
  \item Avoid contradicting its own earlier statements
  \item Converge toward task completion
  \item Use the Focus Economy efficiently
\end{enumerate}

\subsection{Memory Formation Quality}

When granted access to memory tools, does the model:

\begin{enumerate}
  \item Save facts with appropriate granularity (not too broad, not too specific)
  \item Avoid duplicate saves (testing the model's own dedup awareness, independent of
        the Hippocampus's cosine threshold)
  \item Retrieve relevant memories when contextually appropriate
  \item Build on previous memories in subsequent sessions
\end{enumerate}

\subsection{Context Window Utilization}

As the context fills over multiple turns:

\begin{enumerate}
  \item At what fill percentage does quality degrade?
  \item Does the model handle the River of Six addon's age-decayed history correctly?
  \item Can it distinguish fresh context from aged context?
\end{enumerate}

\subsection{Quantization Effects}

For a given model family (e.g., Llama 3.1), compare across quantizations:

\begin{enumerate}
  \item q4\_0 vs. q4\_K\_M vs. q5\_K\_M vs. q8\_0 vs. fp16
  \item At what quantization level do tool calls become unreliable?
  \item Does quantization affect multi-turn consistency differently than single-turn quality?
\end{enumerate}

\section{Harness Architecture}

\subsection{Test Suites as Neural Pathways}

% TODO: How test suites are implemented as CNS pathways
% - Each test case is an effector in a pathway
% - Spike blackboard carries test parameters
% - Results collected via spike status + telemetry
% - Standard Are-Self execution infrastructure (no separate test runner)

\subsection{Evaluation Metrics}

% TODO: Define formal metrics
% - Tool Call Accuracy (TCA): valid calls / total calls
% - Focus Efficiency (FE): task progress / focus consumed
% - Memory Precision (MP): relevant saves / total saves
% - Context Degradation Curve: quality vs. context fill %
% - Turns to Solution (TTS): fewer is better

\subsection{Automated Scoring}

% TODO: How scoring works
% - Deterministic checks (JSON format, tool name existence)
% - LLM-as-judge for quality metrics (using a known-good model)
% - Statistical aggregation across multiple runs
% - Variance analysis for reliability

\section{Evaluation}

% TODO: Results across model families
% - Llama 3.1 (7B, 13B, 70B) across quantizations
% - Gemma 2/4 series
% - Mistral/Mixtral
% - Qwen 2.5 series
% - Phi-3/4 series
% - Focus on 7B models (target hardware constraint)

\section{Discussion}

% TODO: Key findings
% - Which models work best as Are-Self agents?
% - Surprising quantization effects
% - Recommendations for model selection
% - How the Hypothalamus's model routing can use harness results

\section{Conclusion}

We have presented a testing harness that evaluates local LLMs along the dimensions that
matter for autonomous agent operation: tool calling, multi-turn consistency, memory
formation, and context utilization. By implementing the harness within Are-Self's own
execution infrastructure, evaluations run in the exact environment where models will
operate, producing results that directly predict real-world agent performance.

\printbibliography

\end{document}
